% !TeX root = ../../Book3.tex
\OptionalChapter{统一化定理与共形类型}{b3:ch:12}
\BookChapterNumber{b3:ch:12}
\begin{chapterabstract}
本选读章在三种单连通模型上准确陈述深层统一化定理，注明完整证明的正式来源与本章证明范围。随后从甲板群完整推导商曲面的分类和共形类型，解释复环面、穿孔曲面及平面域的抛物型例外，并展示 Poincaré 密度如何由圆盘覆盖下降。
\end{chapterabstract}

\begin{learninggoals}
\begin{itemize}
\item\label{b3:item:12-goal-1}区分单连通曲面的统一化与平面域的 Riemann 映射。
\item\label{b3:item:12-goal-2}用紧性、Liouville 和固定点区分三种模型。
\item\label{b3:item:12-goal-3}从自由甲板群推导平面、去心平面、复环面及圆盘商。
\item\label{b3:item:12-goal-4}准确解释覆盖共形类型及后续主线所需的深层输入。
\end{itemize}
\end{learninggoals}

覆盖理论把曲面的拓扑复杂度提升到一张单连通曲面，但没有给它一个全局复平面坐标。统一化完成这一步，并使拓扑路径类变成标准模型上的全纯自同构群。

\ProseParagraphBreak
本章为选读层次，深层定理的完整证明未声称在此完成。其推论的假设、证明和反例均独立写明；第13章将重新陈述所需覆盖输入，避免选读章成为不可绕开的唯一前提。

% !TeX root = ../../Book3.tex
\section{单连通 Riemann 曲面}
\label{b3:sec:1201}

\chapref{b3:ch:11}已经构造每张连通 Riemann 曲面的全纯万有覆盖，并证明底曲面是源曲面对甲板群的商。现在还缺少一个解析问题：单连通的源曲面能够采用什么全局复坐标？本节先研究球面、平面、圆盘三个具体模型，算清前两者允许的自由甲板群，并比较圆盘的两种循环商。\secref{b3:sec:1202}的统一化定理才负责证明这些模型穷尽全部可能。

\subsection{球面}
\label{b3:subsec:120101}

Riemann 球面 \(\widehat{\mathbb C\kern0pt}\) 使用有限坐标 \(z\) 与无穷远坐标 \(w=1/z\)，重叠处的转换全纯且导数非零。两个坐标闭圆盘 \(|z|\le1\)、\(|w|\le1\) 覆盖全曲面，所以它紧。

\ProseParagraphBreak
球面还单连通，但不能以“每条连续闭路都漏掉一个点”作为理由，连续路径的像可能非常复杂。对一条给定闭路，用有限小坐标圆盘覆盖其像，再细分参数，使每段落在一个坐标圆盘内。在坐标中把该段直线替换，凸性给出端点固定同伦。替换后的路径只由有限条圆弧或线段组成，可以选一个未被经过的球面点 \(q\)。删去 \(q\) 后的球面经一个 Möbius 坐标等价于平面，替换路径便可在其中收缩。原闭路也随之可收缩。这里先作有限局部替换，才使用去掉一点的平面坐标。

\ProseParagraphBreak
球面上的亚纯函数已经由\cref{b3:thm:03-sphere-rational}证明为有理函数。自同构比一般有理函数更严格：它在每个点都有全纯局部逆，因而不能出现重复覆盖。下面用唯一极点及其阶数把这一限制落实。

\Needspace{6\baselineskip}
\begin{proposition}[球面不能有非平凡自由甲板群]\label{b3:prop:12-sphere-free}
每个双全纯自同构 \(T:\widehat{\mathbb C\kern0pt}\to\widehat{\mathbb C\kern0pt}\) 都是 Möbius 变换
\[
 T(z)=\frac{az+b}{cz+d},\quad a,b,c,d\in\mathbb C,
 \quad ad-bc\ne0.
\]
若 \(T\) 非恒等，则它至少有一个球面不动点。因此，若 \(X\) 为连通 Riemann 曲面，\(p:\widehat{\mathbb C\kern0pt}\to X\) 为全纯万有覆盖，则 \(p\) 是双全纯映射。
\end{proposition}
\begin{proof}
令 \(q=T^{-1}(\infty)\)。由于 \(T\) 单射，\(q\) 是它唯一的极点；在源坐标及目标无穷远坐标中，\(T\) 有全纯局部逆，所以这个极点是一阶的。选 Möbius 变换 \(M\)，使 \(M(\infty)=q\)：若 \(q\) 有限可取 \(M(z)=q+1/z\)，否则取恒等映射。于是 \(R=T\circ M\) 只在无穷远有一个一阶极点，在有限平面全纯。

将球面全纯映射看作亚纯函数，\cref{b3:thm:03-sphere-rational}给 \(R=P/Q\)。约去公因子后，分母不能有有限零点，故由代数基本定理为常数；因此 \(R\) 是多项式。它在无穷远的极点阶数就是多项式次数，所以次数为1。\(T=R\circ M^{-1}\) 因而是 Möbius 变换，且行列式非零。

若 \(c=0\)，\(T\) 固定无穷远。若 \(c\ne0\)，有限不动点满足
\[
 cz^2+(d-a)z-b=0.
\]
代数基本定理保证至少一个根。这个根不可能同时使 \(cz+d=0\)，否则原方程还会给 \(az+b=0\)，违背 \(ad-bc\ne0\)。所以它确实是有限不动点。

最后，甲板变换全纯，而\cref{b3:prop:11-deck-free}说明非恒等甲板变换不能固定源点，故 \(p\) 的甲板群只有恒等元。由于 \(p\) 是万有覆盖，\cref{b3:prop:11-universal-deck}保证甲板群在每条纤维上传递，所以每个纤维只有一点。覆盖又局部双全纯，因而这个全局双射及其逆都全纯。
\end{proof}
这里起作用的是“所有球面自同构都有不动点”与“甲板作用自由”的结合。紧性本身不能排除非平凡覆盖：后面的平移格商也紧，却由整个平面覆盖。另一方面，球面上的全局复值全纯函数必恒常，因为模在紧球面上取得最大值，局部最大模原理及连通性迫使函数恒常；有理函数的丰富性来自允许取值 \(\infty\)。

\subsection{平面}
\label{b3:subsec:120102}

平面 \(\mathbb C\) 凸，所以单连通；它没有球面的紧性。\Cref{b3:thm:05-entire-univalent}已经证明，单射整函数恰为 \(z\mapsto az+b\)、\(a\ne0\)。因此平面自同构的自由作用受到一个直接限制：若 \(a\ne1\)，点 \(b/(1-a)\) 被固定；非恒等甲板变换只能是非零平移。

\ProseParagraphBreak
多个平移能够组成什么群，取决于它们能否无限接近恒等变换。覆盖要求每个小层与所有非恒等甲板像不交，因而排除了任意短的非零平移。以下证明把这项局部限制转成平面上的整数格。

\Needspace{6\baselineskip}
\begin{proposition}[平面的离散平移群]\label{b3:prop:12-translation-lattice}
设 \(X\) 为连通 Riemann 曲面，\(p:\mathbb C\to X\) 为全纯万有覆盖。则所有甲板变换均为平移，且平移向量群
\[
 \Lambda=\{\,\lambda\in\mathbb C\mid
              (z\mapsto z+\lambda)\in\operatorname{Deck}(p)\,\}
\]
恰为下列三种形式之一：\(\{0\}\)、\(\mathbb Zu\)、\(\mathbb Zu+\mathbb Zv\)。第二种中 \(u\ne0\)；第三种中 \(u,v\) 在 \(\mathbb R\) 上线性独立。反过来，每个所列加法群的平移作用都自由且适当不连续，其商具有以 \(\mathbb C\) 为源的全纯万有覆盖。
\end{proposition}
\begin{proof}
由平面自同构分类及甲板作用自由性，每个甲板变换均为平移，向量在加法与取负下封闭。由\cref{b3:prop:11-deck-proper}，0有邻域 \(W\)，与其所有非恒等甲板像不交。选 \(\varepsilon>0\)，使 \(D(0,\varepsilon)\subset W\)。若非零 \(\lambda\in\Lambda\) 满足 \(|\lambda|<\varepsilon\)，则 \(\lambda\) 同时属于 \(W\) 与 \(W+\lambda\)，矛盾。因此每个非零向量的模都至少为 \(\varepsilon\)。

更重要的是，任意两个不同向量之差仍是非零群元，所以
\[
 |\lambda-\mu|\ge\varepsilon
 \quad(\lambda,\mu\in\Lambda,\ \lambda\ne\mu).
\]
用有限个半径小于 \(\varepsilon/2\) 的圆盘覆盖任一闭有界圆盘，每个小圆盘至多含一个群元，故 \(\Lambda\) 与任何有界集的交都是有限集。这一步使用了群的差封闭性，而不仅是逐点孤立。

若 \(\Lambda\ne\{0\}\)，先任选非零 \(\lambda_0\)。集合
\(\{\lambda\in\Lambda:0<|\lambda|\le|\lambda_0|\}\)
有限非空，从中取模最小的 \(u\)，便是全群最短非零向量。若 \(tu\in\Lambda\)、\(t\in\mathbb R\)，选整数 \(n\)，使 \(|t-n|\le1/2\)。向量 \((t-n)u\) 仍在群中，且模不超过 \(|u|/2\)，最短性迫使它为零。因此
\[
 \Lambda\cap\mathbb Ru=\mathbb Zu.
\]
若全部向量都在这条实直线上，得到第二种形式。

否则以 \(u\) 归一化坐标，写
\[
 \frac wu=a(w)+\mathrm i b(w),\quad a(w),b(w)\in\mathbb R.
\]
取 \(w_0\in\Lambda\)，经必要的变号使 \(b_0=b(w_0)>0\)。减去整数倍 \(u\) 可以把任何向量的横坐标归约到 \([0,1)\)，而不改变高度 \(b(w)\)。于是集合
\[
 E=\{\,w\in\Lambda\mid0\le a(w)<1,\quad0<b(w)\le b_0\,\}
\]
非空，且包含在一个有界平行四边形中，所以有限。取 \(v\in E\)，使 \(h=b(v)\) 最小。若全群还有正高度小于 \(h\) 的向量，先归约其横坐标，它也会落在 \(E\)，与最小性矛盾。故 \(h\) 是整个群的最小正高度；这同时证明正高度有正下界并且下界取得。

对任意 \(w\in\Lambda\)，取 \(n=\lfloor b(w)/h\rfloor\)，则
\[
 0\le b(w-nv)<h.
\]
最小正高度使这个高度只能为零。因此 \(w-nv\in\Lambda\cap\mathbb Ru=\mathbb Zu\)，得到 \(\Lambda=\mathbb Zu+\mathbb Zv\)。因为 \(b(v)>0\)，两个向量在实数域上线性独立。

反过来，所列群中的整数坐标使它们与有界集只有有限交点；秩二情形使用实线性同构 \((s,t)\mapsto su+tv\)，把有界集拉回有界集。非零平移没有固定点，故作用自由。若紧集 \(K,L\subset\mathbb C\) 满足 \((K+\lambda)\cap L\ne\varnothing\)，则 \(\lambda\in L-K\)，而 \(L-K\) 紧，只有有限个这样的群元。由\cref{b3:prop:11-proper-action-quotient}，商成为 Riemann 曲面，商投影是全纯覆盖；源平面单连通，故为万有覆盖。
\end{proof}
三种形式中的秩，分别指0、1、2个独立的整数生成元。两个实方向已被最短向量与最小高度用尽，不能在保持离散性的同时再加入第三个整数独立方向。商曲面是什么，则还须把坐标写出来。

\ProseParagraphBreak
秩0时没有不同点被识别，商就是平面。秩1时令
\[
 q_u(z)=\exp\left(\frac{2\pi\mathrm i z}{u}\right).
\]
指数的周期说明 \(q_u(z)=q_u(w)\) 当且仅当 \(z-w\in\mathbb Zu\)。每个非零复数都有对数，所以 \(q_u\) 满射到 \(\mathbb C^*\)；其导数处处非零。因此它在 \(\mathbb C/\mathbb Zu\) 上诱导一个双射，且在每个小商坐标中双全纯，得到
\[
 \mathbb C/\mathbb Zu\simeq\mathbb C^*.
\]
这项结论包含全部纤维和局部逆的核验，而不只是写出一个周期函数。

\ProseParagraphBreak
秩2时，每个 \(z\in\mathbb C\) 唯一写成 \(z=su+tv\)、\(s,t\in\mathbb R\)。分别减去 \(s,t\) 的整数部分，表明闭平行四边形
\[
 P=\{\,su+tv\mid0\le s\le1,\ 0\le t\le1\,\}
\]
映满商曲面，故商紧。上下两边及左右两边分别按平移识别。更明确地，记 \(\mathbb S^1=\{\zeta:|\zeta|=1\}\)，映射
\[
 [su+tv]\longmapsto
 \bigl(e^{2\pi\mathrm i s},e^{2\pi\mathrm i t}\bigr)
\]
是商到 \(\mathbb S^1\times\mathbb S^1\) 的连续双射；源紧而目标 Hausdorff，所以它是同胚。这给出了环面的拓扑，而复结构由平面小圆盘的商坐标给出。这样的 Riemann 曲面称为\term[fu-huanmian]{复环面}[complex torus]。实坐标表示说明其拓扑，并未把任意实线性变换宣布为全纯坐标变换。

\subsection{圆盘}
\label{b3:subsec:120103}

圆盘 \(\mathbb D\) 也凸、单连通且非紧。它与平面有相同的开圆盘拓扑，却有不同的全纯函数约束：恒等函数是非恒常有界全纯函数，而平面的有界整函数由 Liouville 定理只能常值。圆盘的自同构则通过 Cayley 变换成为上半平面的实 Möbius 变换。

\ProseParagraphBreak
由\cref{b3:prop:11-mobius-types}，一个没有内部固定点的非恒等圆盘自同构，在半平面坐标中属于抛物型或双曲型。对它生成的循环群，必要时更换生成元或缩放坐标，可以分别使用
\[
 T(w)=w+1,\quad S_a(w)=aw,\quad a>1
\]
作为模型。\subsecref{b3:subsec:110401}已经给出全部纤维及局部逆的计算，所得商为
\[
 \begin{aligned}
 \mathbb H/\langle T\rangle
 &\simeq\mathbb D\setminus\{0\},\\
 \mathbb H/\langle S_a\rangle
 &\simeq A\left(e^{-2\pi^2/\log a},1\right).
 \end{aligned}
\]
前者的商坐标是 \(e^{2\pi\mathrm i w}\)；后者在 \(0<\operatorname{Im}\log w<\pi\) 的分支中使用 \(\exp(2\pi\mathrm i\log w/\log a)\)。这些映射分别把整周期与伸缩轨道识别成同一点。

\ProseParagraphBreak
两商的基本群都为 \(\mathbb Z\)：它们的源单连通，甲板群又是无限循环群，可用\cref{b3:prop:11-fundamental-deck-action}识别。它们的万有覆盖也都等价于圆盘，却仍有不同的共形结构。双曲型循环商的模量为
\[
 \frac1{2\pi}\log\left(e^{2\pi^2/\log a}\right)
 =\frac{\pi}{\log a}.
\]
例如 \(a=e^\pi\) 给模量1，\(a=e^{2\pi}\) 给模量 \(1/2\)，由\cref{b3:prop:08-modulus-invariance}，这两个圆环不双全纯等价。去心盘的两端连接族具有无穷极值长度，\secref{b3:sec:0803}已用趋向穿孔的子圆环密度证明这一点；所以它也不等价于任何有限模量圆环。

\ProseParagraphBreak
平面的秩1商 \(\mathbb C^*\) 同样有基本群 \(\mathbb Z\)，但也不能与去心盘等价。否则将一个到去心盘的双全纯映射与指数 \(\mathbb C\to\mathbb C^*\) 复合，便得到非恒常有界整函数，违反 Liouville。由此，抽象基本群既不能决定万有覆盖的共形模型，也不能决定给定模型上甲板群的具体作用。

\subsection{三种基本模型}
\label{b3:subsec:120104}

以上讨论从具体模型出发，尚未把任意单连通曲面放进其中。关键区别在于全局坐标是否已经存在：Riemann 曲面的定义只给一组相容的局部坐标，Riemann 映照定理则以一个已经位于平面中的区域为出发点。下面把可以直接使用该定理的范围单独说清。

\Needspace{6\baselineskip}
\begin{proposition}[已嵌入球面的单连通曲面]\label{b3:prop:12-embedded-models}
设 \(X\) 为连通、单连通 Riemann 曲面。若存在全纯单射 \(F:X\to\widehat{\mathbb C\kern0pt}\)，则 \(X\) 双全纯等价于 \(\widehat{\mathbb C\kern0pt}\)、\(\mathbb C\)、\(\mathbb D\) 中恰好一种。
\end{proposition}
\begin{proof}
在源与目标的局部坐标中，单射全纯函数由\cref{b3:cor:01-local-injectivity}导数非零，所以 \(F\) 在每点局部双全纯。它的像 \(U\) 开，且局部逆与单射性拼成 \(F:X\to U\) 的全纯逆，故 \(U\) 单连通。

若 \(U\) 为整个球面，结论已得。否则取 \(a\notin U\)，用 Möbius 变换把 \(a\) 送到无穷远。像于是成为一个单连通平面区域 \(V\)。若 \(V=\mathbb C\)，得到平面模型；若 \(V\subsetneq\mathbb C\)，\cref{b3:thm:05-riemann}给 \(V\) 到圆盘的双全纯映射。

最后核对唯一性。球面紧，平面与圆盘不紧，故球面不能与后二者双全纯。若平面与圆盘双全纯，就存在非恒常有界整函数，与\cref{b3:thm:02-liouville}矛盾。因此三模型两两不等价。
\end{proof}
这项命题已完成“有全局单射坐标以后”的分类，却没有产生这个坐标。即使找到一个非恒常有界全纯函数，也还须证明它单射；圆盘上的 \(z\mapsto z^2\) 就表明有界性与非常值性不够。对紧单连通曲面，全局坐标还必须允许取无穷远值，不能只寻找复值全纯函数。

\ProseParagraphBreak
统一化的深层内容，正是对任意抽象单连通 Riemann 曲面构造这样的全局模型，而不是把已有平面区域再分类一次。有了覆盖模型，球面不动点、离散平移群与循环商的计算就能帮助辨认底曲面的形状。

% !TeX root = ../../Book3.tex
\section{统一化定理}
\label{b3:sec:1202}

统一化将复杂曲面的局部复结构汇成一个全局模型。本书在这里准确陈述完整深层结论并说明证明路线；后续关于群商、密度与全纯映射的推论另作完整证明。

\subsection{定理陈述}
\label{b3:subsec:120201}

\term[tongyihua-dingli]{统一化定理}[uniformization theorem]识别所有单连通曲面的复结构模型。

\begin{theorem}[深层统一化定理]\label{b3:thm:12-uniformization}
每张连通单连通 Riemann 曲面双全纯等价于 \(\widehat{\mathbb C\kern0pt}\)、\(\mathbb C\)、\(\mathbb D\) 三者之一，且模型唯一。
\end{theorem}
\begin{proofsketch}
一般曲面统一化的完整证明不在本章重建。已核查的正式来源是 Ahlfors《Conformal Invariants》\cite{Ahlfors2010Conformal}第10章，定理10-3（142页），证明在 \S10-4 与 \S10-5（142--149页）。有 Green 函数的单连通情形用其调和共轭构造有界全纯坐标，再证明该坐标单叶且满圆盘；不存在 Green 函数时，利用调和比较和极值构造带一个指定简单极点的亚纯函数，证明它全局单叶，从而得到平面模型；紧情形的类似构造给球面。Green 函数的构造、极值函数的全局单射性和抛物型论证均为必须补足的深层环节，本段只说明路线。三模型唯一性已经由前节完整证明。
\end{proofsketch}
这里的引用明确区分“采用定理”与“重建证明”。特别是无 Green 函数不能单靠一句 Liouville 反向推出平面坐标存在。

\subsection{单连通形式}
\label{b3:subsec:120202}

统一化不依赖预先给一条光滑边界，也不要求曲面可嵌入三维欧氏空间。曲面的坐标变换只要全纯相容，深层定理就提供某个标准模型。这个模型中的坐标拉回原曲面后，才是全局单值参数。

\ProseParagraphBreak
在平面单连通区域上，定理退化为已证明的 Riemann 映射及全平面的特例；在抽象曲面上，它增加了全局嵌入这一新内容。学习时可先掌握这些特殊情形，再阅读深层证明的势论路线。后面的主线章会独立列出所需的覆盖类型输入，因此本选读章不是它们唯一的逻辑前置。

\ProseParagraphBreak
第五章的 Riemann 映射定理是平面区域的定理；本章的统一化允许源曲面没有全局平面坐标。不能先把一般单连通曲面“展开为一个平面区域”再套第五章，因为这种展开本身就是需要证明的全局结论。局部坐标图在重叠处全纯，只保证复结构相容，并不保证全曲面可单射放入平面。球面甚至因紧性不可能双全纯嵌到平面区域。

\ProseParagraphBreak
深层证明通常把曲面上的调和函数或 Green 函数作为构造工具，分别处理存在 Green 函数与不存在的情形，再建立模型映射。单连通性负责消除共轭调和函数的周期，极值或穷竭负责全局单射与满射。这里只解释这些步骤的任务，未给曲面势论的完整存在理论；所以定理明确标为深层输入，并用正式证明来源承接。后面的度量公式则可在此输入以后逐项核验，而无须重复整个存在证明。

\subsection{万有覆盖形式}
\label{b3:subsec:120203}

\begin{corollary}\label{b3:cor:12-quotient-uniformization}
每张 Riemann 曲面 \(X\) 可写为三模型之一对一个自由适当不连续全纯自同构群的商。模型是其万有覆盖，群是甲板群。
\end{corollary}
\begin{proof}
第11章构造了单连通全纯万有覆盖，并证明底曲面等于甲板群商。对覆盖曲面使用统一化，选一个双全纯模型坐标，把甲板群共轭到模型上；自由性与局部不交性质在共轭下保持，所以得到所述商。
\end{proof}
若模型平面，群只能由平移组成。若模型球面，群只能平凡。圆盘模型则允许许多非交换离散群，所以高拓扑复杂度集中于这一情形。

\subsection{共形类型的唯一性}
\label{b3:subsec:120204}

\begin{proposition}\label{b3:prop:12-cover-type-unique}
同一曲面的两个单连通全纯万有覆盖双全纯同构，因此覆盖模型与所选覆盖无关。
\end{proposition}
\begin{proof}
指定同一底点的两个纤维点。沿底路径提升得到唯一覆盖同构，其在每个小层的坐标中为恒等，故双全纯。若两模型不同，这个同构会给不同模型之间的双全纯映射，违反前节区分。因此模型唯一。
\end{proof}
不同的模型坐标只把甲板群共轭。甲板群写成具体矩阵时可能变化，群商的曲面及共形类型并不变化。

\ProseParagraphBreak
三模型之间不能双全纯等价，已有工具分别排除。球面紧，平面与圆盘不紧；若平面与圆盘双全纯，一张整个平面的有界非恒定全纯函数就出现，违反 Liouville。唯一性指这三种模型中的类型唯一，不是说圆盘坐标唯一：圆盘自同构仍给无限多坐标选择，平面有仿射自由度，球面有 Möbius 自由度。通常以一个点及导数归一化，或三个球面点归一化，来消去相应自由度。

\ProseParagraphBreak
万有覆盖也只在选好基点提升后唯一。两张覆盖可由唯一保持所选基点的同构对应；改变提升点，差一个甲板变换。这些变换恰保存后面的 Poincaré 度量，所以度量不会依赖覆盖坐标的选择。若没有先证明甲板不变性，仅说“选任意覆盖映射定义密度”，良定义仍未完成。

% !TeX root = ../../Book3.tex
\section{双曲型、抛物型与椭圆型}
\label{b3:sec:1203}

共形类型指万有覆盖的模型，不是底曲面是否有界、是否紧或是否有孔的简单同义词。特别是紧曲面也可能有圆盘万有覆盖。

\subsection{共形类型}
\label{b3:subsec:120301}

\begin{definition}\label{b3:def:12-conformal-types}
万有覆盖分别为球面、平面、圆盘的曲面称为\term[tuoyuanxing-qumian]{椭圆型曲面}[elliptic Riemann surface]、\term[paowuxing-qumian]{抛物型曲面}[parabolic Riemann surface]、\term[shuangquxing-qumian]{双曲型曲面}[hyperbolic Riemann surface]。这三个词在本卷始终按覆盖模型解释。
\end{definition}
势论中的“有无 Green 函数”也常称为双曲与抛物分类，但对一般非单连通曲面，它与这里的覆盖分类不宜无条件等同。紧的高亏格曲面具有圆盘覆盖，却没有单个对数极点的全局 Green 函数；紧性使 Laplace 方程需要额外的零平均补偿。

\ProseParagraphBreak
本书用“覆盖双曲型”或明确提到圆盘覆盖消除这个语境差别。不要只由一个词的英文相同就把不同分类中的性质互相搬用。

\ProseParagraphBreak
曲面类型说的是万有覆盖，并不要求曲面本身等于那个模型。有限环域、去心盘及紧亏格至少2曲面都为圆盘型，却分别有两条普通边界、一个穿孔端或没有边界的紧拓扑。平面型也包含去心平面和紧环面，而非只包含平面。术语若用在一般曲面，必须把覆盖定义保留；对于单连通曲面才可直接把它与模型等同。

\ProseParagraphBreak
尤其不能把本章的抛物型直接读成“没有 Green 函数”的全部一般曲面分类。对单连通非紧模型两种观点相互对应，但对一般商曲面，势论是否存在 Green 函数还与商的全局结构有关。本章只按覆盖定义命名，后面的势论会按其自己的存在条件陈述。相同形容词在不同分类框架中的定义范围必须看清。

\subsection{万有覆盖}
\label{b3:subsec:120302}

平面甲板群对应一个离散加法子群 \(\Lambda\subset\mathbb C\)。局部不交性给原点附近没有非零群元。若没有非零元，商为平面；若群位于一条实直线上，最短正方向元生成它，商为去心平面；若存在两个实线性无关元，取其格的基本平行四边形，离散性使其中只有有限候选，再缩减可选出一个格基，群为 \(\mathbb Z\omega_1+\mathbb Z\omega_2\)。其商为复环面。

\ProseParagraphBreak
因此抛物型曲面恰有平面、去心平面、复环面这三类形式；复环面仍有不同模参数，绝非只有一个共形类。球面型只有球面。其他曲面的覆盖模型只能为圆盘。这个分类由自由甲板变换与离散平移群得到，不需要假定底曲面有一个平面单值坐标。

\subsection{例子}
\label{b3:subsec:120303}

标准有限圆环的万有覆盖是带，带共形等价于圆盘，所以有限圆环双曲型。去心圆盘覆盖为半平面，也双曲型；去心平面覆盖为全平面，因此抛物型。三者基本群都同构 \(\mathbb Z\)，故基本群的抽象同构类不能区分它们的共形类型。

\ProseParagraphBreak
复环面 \(\mathbb C/\Lambda\) 紧而抛物型；球面紧而椭圆型；亏格至少2的紧曲面双曲型。最后一句可由其基本群非交换排除平移商，也由非球面拓扑排除球面。底曲面紧性只决定商是否紧，并不决定万有覆盖是否紧。

\subsection{穿孔与紧曲面}
\label{b3:subsec:120304}

从球面删去一点得平面，删去两点得去心平面，删去至少三点则为双曲型平面区域。这展示了穿孔数量改变覆盖模型的临界现象。对复环面删去一点，得到非紧且基本群非交换的曲面，也双曲型。

\ProseParagraphBreak
“椭圆型”与“椭圆函数”不是同一个分类。椭圆函数生活在平移格商复环面上，其万有覆盖是平面，按本卷定义为抛物型。名称的历史来源不同，不能据字面把环面放入球面型。类似地，“双曲型”也不表示曲面作为欧氏图形有双曲线边界。

\ProseParagraphBreak
穿孔可能改变类型。球面删去一点是平面，删去两点是 \(\mathbb C^*\)，删去三点已为圆盘型。三点的位置可用 Möbius 变换归一化为 \(0,1,\infty\)，得到 \(\mathbb C\setminus\{0,1\}\)。基本群含两条独立绕圈，并非平面型商的秩0或秩1平移群；同时区域不紧，不能是秩2平移商的环面。这与一般平面双曲分类的证明相符。

\ProseParagraphBreak
环面删去一点也变成圆盘型。若仍是平面型，上述分类只允许平面、去心平面或完整环面；前三者的拓扑分别为无孔、一个基本圈与紧面，不等于去掉一点的环面。这里使用曲面的通常拓扑分类或基本群计算。没有把局部删去一个圆盘的图误当作整个曲面；穿孔附近看起来都像去心盘，但全局剩余的柄仍影响覆盖类型。

\ProseParagraphBreak
亏格至少2的紧曲面不能是球面型，也不能是平面型：前者只有球面，后者的唯一紧商是环面。因而由统一化它必为圆盘型。亏格1的紧曲面反向确为环面，需要曲面的全局函数或1形式理论；附录A从 Riemann--Roch 给出这一联系。若只知道一个给定商 \(\mathbb C/\Lambda\)，其环面结构直接可见，不需要反向分类。区分“所构造模型属于某类型”与“所有这种拓扑的复结构都属于某类型”，可避免在例子中遗漏真正的存在论证。

% !TeX root = ../../Book3.tex
\section{平面域的统一化}
\label{b3:sec:1204}

一般曲面的抛物型包括复环面，平面区域却不能容纳这个紧例外。还需回答一个看似较小的问题：一张与 \(\mathbb C^*\) 双全纯的区域，是否可能在平面中缺少许多点？不能仅凭它有一个孔就作判断；下面利用单叶函数在两端的奇点，把实际像确定下来。

\subsection{平面双曲区域}
\label{b3:subsec:120401}

\Needspace{6\baselineskip}
\begin{lemma}[去心平面的单叶平面像]\label{b3:lem:12-punctured-plane-embedding}
设 \(f:\mathbb C^*\to\mathbb C\) 为单射全纯函数。则存在 \(A\in\mathbb C^*\)、\(B\in\mathbb C\)，使
\[
 f(z)=Az+B\quad\hbox{或}\quad f(z)=\frac A z+B.
\]
特别地，\(f(\mathbb C^*)=\mathbb C\setminus\{B\}\)。
\end{lemma}
\begin{proof}
先排除0处的本性奇点。取 \(z_0\in\mathbb C^*\)，再取闭包避开0的一个小圆盘 \(V\ni z_0\)。开映射定理使 \(f(V)\) 包含 \(f(z_0)\) 的某个邻域。若0为本性奇点，Casorati--Weierstrass 定理给 \(z_n\to0\)，使 \(f(z_n)\to f(z_0)\)。充分大的 \(n\) 满足 \(z_n\notin V\)、\(f(z_n)\in f(V)\)，于是这个值既有 \(V\) 中的原像，又有 \(z_n\) 这个原像，违反单射。对 \(f(1/w)\) 在 \(w=0\) 重复相同论证，也排除了无穷远处的本性奇点。

因此两端都是可去奇点或极点。补入两端后得到球面亚纯函数，由\cref{b3:thm:03-sphere-rational}，它是有理函数 \(R=P/Q\)，其中 \(P,Q\) 互素。令 \(d=\max(\deg P,\deg Q)\)。可选一个有限值 \(c\)，避开 \(R(0)\)、\(R(\infty)\)、所有临界值，以及使 \(P-cQ\) 的最高次项消掉的至多一个值。方程 \(P-cQ=0\) 于是有 \(d\) 个互异的有限根；互素性保证这些根不是极点，而所选 \(c\) 保证它们均不为0。每个根都是原函数 \(f\) 的原像，单射性迫使 \(d=1\)。

所以 \(R\) 为 Möbius 变换。它在球面恰有一个极点，该点不能落在 \(\mathbb C^*\)，故只能是0或无穷远。极点为无穷远时 \(R=Az+B\)；极点为0时 \(R=A/z+B\)，且 \(A\ne0\)。两式的像直接给出结论。
\end{proof}
这里单射性既排除了本性奇点，也排除了高次有理映射。若只假设导数不为零，\(e^z\) 在无穷远仍有本性奇点；局部可逆不足以完成这一步。

\Needspace{6\baselineskip}
\begin{proposition}[平面区域的圆盘覆盖]\label{b3:prop:12-planar-hyperbolic}
采用深层统一化定理。设 \(U\subseteq\mathbb C\) 是非空连通开集。则 \(U\) 具有全纯圆盘万有覆盖，当且仅当 \(\mathbb C\setminus U\) 至少含两个不同点。其余情形恰为 \(U=\mathbb C\) 或 \(U=\mathbb C\setminus\{a\}\)，二者都具有平面万有覆盖。
\end{proposition}
\begin{proof}
先设补集至少有两个点。统一化给 \(U\) 的三种可能覆盖模型。球面模型的商只有球面，而非空平面区域不紧，故排除球面。若为平面模型，由前节的平面商分类，\(U\) 必与 \(\mathbb C\)、\(\mathbb C^*\) 或某个复环面双全纯。紧复环面仍被排除。

若存在双全纯映射 \(f:\mathbb C\to U\)，把它视为单叶整函数，\cref{b3:thm:05-entire-univalent}给 \(f(z)=Az+B\)，从而 \(U=\mathbb C\)。若存在双全纯映射 \(f:\mathbb C^*\to U\)，前一引理给 \(U=\mathbb C\setminus\{B\}\)。两种都与至少缺少两个点矛盾，因此只剩圆盘模型。

反向，补集为空时，恒等映射就是平面万有覆盖；补集只有 \(a\) 时，映射 \(w\mapsto a+e^w\) 是全纯覆盖，源平面单连通，故为平面万有覆盖。模型唯一性排除这两种区域再有圆盘万有覆盖。
\end{proof}
这是平面补集的计数；换到球面，还要把无穷远算作一个遗漏点。因此球面域的相应门槛是至少遗漏三个点。证明不要求补集有限，也不要求边界光滑或局部连通。只要选出两个不同的遗漏点，剩下的补集可以具有任意复杂形状。

\ProseParagraphBreak
内部圆盘覆盖与边界对应仍是不同结论。这里的投影在每个内部点附近都有完整的局部层，却不保证延伸为闭圆盘到普通闭包的连续映射。第七章关于边界与素端的条件仍须另行核对。

\subsection{特殊抛物型情形}
\label{b3:subsec:120402}

“真平面域”足以在单连通假设下使用 Riemann 映照，却不足以在任意区域上保证圆盘覆盖。\(\mathbb C^*\) 就是真平面域，它的指数覆盖是整个平面。与之对照，去心盘 \(\mathbb D^*\) 有上半平面覆盖，而有限圆环有条带覆盖。三者都是一个圆圈的同伦类型，复结构却区分了它们的端部。

\ProseParagraphBreak
删点的临界现象也可在球面上直接列出。不给球面删点，模型仍为球面；删去一点，用把该点送到无穷远的 Möbius 坐标得到平面；删去两点，可把它们送到 \(0,\infty\)，得到去心平面。删去至少三个点时，先把一个遗漏点送到无穷远，余下平面补集至少有两个点，前一命题遂给圆盘覆盖。

\ProseParagraphBreak
特别地，三穿孔球面可归一化为 \(\mathbb C\setminus\{0,1\}\)。统一化保证其圆盘覆盖存在，本章并未求出投影的显式公式。许多应用只需要源为圆盘、投影满射且局部双全纯；在这些应用中，缺少公式并不妨碍使用覆盖提升定理。下文的函数刚性就是一个例子。

\subsection{Poincaré 度量构造}
\label{b3:subsec:120403}

圆盘模型还带来一个自然的局部长度单位。由 Schwarz--Pick 的自同构等号，对任意圆盘自同构 \(T\)，
\[
 \frac{2|T'(w)|}{1-|T(w)|^2}=\frac2{1-|w|^2}.
\]
因而密度 \(\lambda_{\mathbb D}(w)\coloneqq2/(1-|w|^2)\) 乘上微小线段长度后，在圆盘坐标变换下不变。要把它传到底曲面，必须先检查两个选择：同一点的不同提升，以及底曲面的不同局部坐标。

\Needspace{6\baselineskip}
\begin{proposition}[圆盘长度元的下降]\label{b3:prop:12-poincare-descent}
设 \(X\) 为连通 Riemann 曲面，\(p:\mathbb D\to X\) 为全纯万有覆盖。在 \(X\) 的局部复坐标 \(z:V\to z(V)\subseteq\mathbb C\) 中，对 \(x\in V\) 任取 \(w\in p^{-1}(x)\)，定义
\[
 \lambda_z(z(x))=
 \frac{2}{(1-|w|^2)\,\bigl|(z\circ p)'(w)\bigr|}.
\]
这里导数在 \(w\) 附近映入 \(V\) 的覆盖小片上计算。这个正光滑密度与提升点、万有覆盖的选择无关；若另一坐标为 \(\zeta=\psi(z)\)，则
\[
 \lambda_\zeta(\psi(z))\,|\psi'(z)|=\lambda_z(z).
\]
因此各坐标中的 \(\lambda_z(z)|\mathrm dz|\) 拼成 \(X\) 上唯一使 \(p\) 局部保持圆盘长度元的共形长度元。
\end{proposition}
\begin{proof}
覆盖局部双全纯，故分母中的导数不为零。若另一个提升为 \(w_2\)，万有覆盖的甲板群在纤维上传递，所以存在圆盘自同构 \(T\)，满足 \(w_2=T(w)\)、\(p\circ T=p\)。在局部坐标中微分得
\[
 (z\circ p)'(w)=(z\circ p)'(T(w))\,T'(w).
\]
结合圆盘密度不变性，得到
\[
 \frac{\lambda_{\mathbb D}(T(w))}
      {|(z\circ p)'(T(w))|}
 =\frac{\lambda_{\mathbb D}(w)}
       {|(z\circ p)'(w)|}.
\]
于是密度与提升无关。固定一个局部逆 \(s:V\to\mathbb D\)，令 \(\sigma=s\circ z^{-1}\)，以 \(u\in z(V)\) 表示坐标数值，则同一公式写成
\[
 \lambda_z(u)=\frac{2|\sigma'(u)|}{1-|\sigma(u)|^2}.
\]
局部逆全纯且导数不为零，故所得密度光滑且正。

换成坐标 \(\zeta=\psi(z)\) 时，链式法则在分母多出 \(|\psi'(z)|\)，恰得陈述中的变换式，因而长度元相同。若改用另一个圆盘万有覆盖 \(p_1\)，覆盖唯一性给圆盘自同构 \(A\)，使 \(p_1\circ A=p\)。以刚才相同的链式法则及密度不变性比较两式，结果仍相同。最后，要求投影局部保持圆盘长度元，就必有
\[
 \lambda_z(z(p(w)))\,|(z\circ p)'(w)|=\lambda_{\mathbb D}(w),
\]
这逐点确定了密度，所以也给唯一性。
\end{proof}
该长度元称为 \term[poincare-duliang]{Poincaré 度量}[Poincaré metric]。它先定义分段光滑曲线的长度
\[
 L(\gamma)=\int \lambda_z(z(\gamma(t)))\,
                 |(z\circ\gamma)'(t)|\,\mathrm dt;
\]
曲线跨过坐标图时分段积分，坐标变换式保证结果一致。路径距离还需对连接两点的曲线长度取下确界。距离公式、完备性和曲率在第十三章独立展开，这里只完成下降本身。

\ProseParagraphBreak
抽象公式可以在两个覆盖上直接算出。先用 Cayley 变换 \(C(w)=(w-\mathrm i)/(w+\mathrm i)\) 把上半平面映到圆盘，有
\[
 1-|C(w)|^2=\frac{4\operatorname{Im}w}{|w+\mathrm i|^2},
 \quad |C'(w)|=\frac2{|w+\mathrm i|^2}.
\]
拉回圆盘长度元，便得上半平面密度 \(1/\operatorname{Im}w\)。这个系数与圆盘中的2对应，后面不再另作缩放。

\Needspace{6\baselineskip}
\begin{proposition}[去心盘与圆环的密度]\label{b3:prop:12-puncture-annulus-density}
设 \(\mathbb D^*=\{z:0<|z|<1\}\)，\(0<r<1\)，\(A(r,1)=\{z:r<|z|<1\}\)，并令 \(L=\log(1/r)\)。以圆盘密度 \(2/(1-|w|^2)\) 归一化的 Poincaré 密度分别为
\[
 \lambda_{\mathbb D^*}(z)=\frac1{|z|\log(1/|z|)},\quad
 \lambda_{A(r,1)}(z)=
 \frac{\pi}{L|z|\sin\!\left(\pi\log(1/|z|)/L\right)}.
\]
\end{proposition}
\begin{proof}
覆盖 \(p(w)=e^{\mathrm i w}\) 将上半平面映到去心盘，且
\[
 |p'(w)|=|p(w)|,\quad
 \operatorname{Im}w=\log(1/|p(w)|).
\]
将上半平面密度除以 \(|p'(w)|\)，得到第一式。

对圆环，使用条带 \(S=\{w:0<\operatorname{Im}w<L\}\) 及同一个公式 \(p(w)=e^{\mathrm i w}\)。映射
\[
 h(w)=e^{\pi w/L}
\]
把 \(S\) 双全纯映到上半平面：辐角在 \((0,\pi)\) 中，模与辐角分别确定实部和虚部。记 \(y=\operatorname{Im}w\)，则
\[
 \frac{|h'(w)|}{\operatorname{Im}h(w)}
 =\frac{\pi}{L\sin(\pi y/L)}.
\]
这就是条带密度。再除以 \(|p'(w)|=|z|\)，并以 \(y=\log(1/|z|)\) 代入，得到第二式。两个投影的纤维均为水平 \(2\pi\mathbb Z\) 平移轨道，条带和半平面都单连通，所以确实使用了全纯万有覆盖。
\end{proof}
密度只依赖 \(|z|\)，但并非常数；靠近不同端部的增长率也不同。去心盘径向曲线从半径 \(\rho_0\) 走到 \(\varepsilon<\rho_0\) 的长度为
\[
 \int_\varepsilon^{\rho_0}\frac{\mathrm d\rho}{\rho\log(1/\rho)}
 =\log\frac{\log(1/\varepsilon)}{\log(1/\rho_0)},
\]
随 \(\varepsilon\downarrow0\) 趋于无穷。这项计算展示穿孔在新长度中的位置；一般路径的距离与完备性还须由后章的完整论证处理。

\subsection{后续章节所需的统一化结论}
\label{b3:subsec:120404}

统一化的一个用途是把全纯映射提升到熟悉的函数空间。目标即使没有有界平面坐标，覆盖坐标也可能把映射化为有界整函数。

\Needspace{6\baselineskip}
\begin{proposition}[平面覆盖源到圆盘覆盖目标]\label{b3:prop:12-parabolic-to-hyperbolic}
设 \(X,Y\) 为连通 Riemann 曲面，\(q:\mathbb C\to X\)、\(p:\mathbb D\to Y\) 分别为全纯万有覆盖。则每个全纯映射 \(f:X\to Y\) 都恒常。
\end{proposition}
\begin{proof}
复合 \(f\circ q:\mathbb C\to Y\) 的源单连通。指定一个源点及其目标提升后，映射提升判别给
\[
 F:\mathbb C\to\mathbb D,\quad p\circ F=f\circ q.
\]
在覆盖小片中 \(F\) 是 \(p\) 的全纯局部逆与 \(f\circ q\) 的复合，因此为有界整函数。Liouville 定理使 \(F\) 恒常，故 \(f\circ q\) 恒常；\(q\) 满射，遂有 \(f\) 恒常。
\end{proof}
例如，一个整函数若遗漏两个不同有限值 \(a,b\)，就给从平面到 \(\mathbb C\setminus\{a,b\}\) 的全纯映射；后者由本节第一项命题有圆盘覆盖，故该整函数恒常。第十四章会把这一路线作为小 Picard 定理重新陈述，并与曲率方法比较。只遗漏一个值时，\(e^z\) 已是非常值反例，恰好落在平面覆盖目标 \(\mathbb C^*\) 上。

\ProseParagraphBreak
后续主线使用的输入可以准确分开。第十三章需要圆盘型曲面的全纯万有覆盖及甲板自同构不变性，然后由圆盘计算和路径提升建立距离理论；第十四章的上述省略值论证只需三穿孔球面的圆盘覆盖存在，不需要其显式模函数。选读章可以略过，但这些章会各自重述所需输入及来源。接受一般统一化、证明其具体后果以及计算一个模型，是三个可以分别检查的步骤。


\begin{chaptersummary}
三种模型按紧性和有界全纯函数区分，万有覆盖模型决定本卷的共形类型。自由平面甲板群只能平移，给平面、去心平面与环面；其他曲面由圆盘商描述。平面补集至少两点保证圆盘覆盖，是后续双曲几何与 Picard 理论的输入。
\end{chaptersummary}

\BookExercises

% 本章题组自拟；以正文已建立的工具设计计算、反例、构造与综合训练。
\begin{exercise}\label{b3:ex:12-1}
说明圆盘和平面拓扑相同而不共形等价。
\end{exercise}
\begin{solution}
径向连续变换可给拓扑同胚。若存在平面到圆盘双全纯映射，它是有界整函数，Liouville 使其常值，矛盾。拓扑同胚不要求全纯。
\end{solution}

\begin{exercise}\label{b3:ex:12-2}
证明球面不存在非平凡自由全纯自同构。
\end{exercise}
\begin{solution}
球面自同构为 Möbius 变换，固定点方程在球面总有解。非恒等变换因此有固定点，不能自由作用。
\end{solution}

\begin{exercise}\label{b3:ex:12-3}
求有限环域、去心盘、去心平面的万有覆盖类型。
\end{exercise}
\begin{solution}
依次为带、半平面、平面。带及半平面与圆盘双全纯，故前两者双曲型，后一者抛物型。三者基本群均为整数群。
\end{solution}

\begin{exercise}\label{b3:ex:12-4}
为何复环面按本卷定义是抛物型而不是椭圆型？
\end{exercise}
\begin{solution}
其万有覆盖是平面，甲板群为格平移；定义按覆盖模型分类。球面覆盖才称椭圆型，椭圆函数的历史名称不决定该分类。
\end{solution}

\begin{exercise}\label{b3:ex:12-5}
给出抛物型曲面中紧与非紧的例子。
\end{exercise}
\begin{solution}
复环面为紧例，平面与去心平面为非紧例。都具有平面万有覆盖，故紧性不能单独判定共形类型。
\end{solution}

\begin{exercise}\label{b3:ex:12-6}
证明平面删去两个不同点具有圆盘万有覆盖。
\end{exercise}
\begin{solution}
补集至少两点，平面域覆盖分类排除球面及平移群的平面、去心平面、紧环面商，只剩圆盘。该结论采用已声明统一化输入。
\end{solution}

\begin{exercise}\label{b3:ex:12-7}
若更换圆盘万有覆盖的模型坐标，解释 Poincaré 密度为何保持相同。
\end{exercise}
\begin{solution}
新坐标与旧坐标由圆盘自同构 \(T\) 联系。恒等式 \(|T'|/(1-|T|^2)=1/(1-|w|^2)\) 使覆盖导数的链式因子恰好抵消，底密度不变。
\end{solution}

\begin{exercise}\label{b3:ex:12-8}
指出“没有非恒常有界全纯函数，所以曲面为平面”的陈述缺少什么。
\end{exercise}
\begin{solution}
必须指定单连通性并采用相应统一化或势论构造。一般复环面也没有非恒常全局复值全纯函数，却是紧曲面而不是平面；否定有界函数不能直接提供一个全局单射坐标。
\end{solution}
